\chapter{La figura del CISO nell'adempimento della Direttiva NIS2}
\label{cap:ciso}

% 5.1 %%%%%%%%%%%%%%%%%%%%%%%%%%%%%%%%%%%%%%%%%%%%%%%%%%%%%%%%%%%%%%%%%%%%%%%%
\section{CISO: compiti ed evoluzione della figura}

\subsection{Compiti del CISO}
Grazie al contributo di Matteo Herin, CISO \& Head of AI Cybersecurity presso Carrefour
Italia, e di Michele Rivieri, CISO di Prometeia, è possibile delineare i principali 
compiti del Chief Information Security Officer in un'organizzazione. A supporto di quanto
reperito dalle interviste, le informazioni sono integrate con il testo \cite{kouns-ciso}.

Come racconta Herin, la funzione cybersecurity nel mondo aziendale ricopre il ruolo
di organo di controllo di primo livello: ha un ruolo di mediazione tra gli obiettivi e 
le scadenze della produzione e la necessità di ottemperare a determinati standard che forniscano
sicurezza all'organizzazione e a tutti i suoi asset, tra cui anche a tutte le persone da
cui è formata.

In quest'ottica, i compiti del CISO si articolano su cinque macrocategorie:
\begin{enumerate}
    \item \textbf{Governance e gestione del rischio}
    \item \textbf{Gestione della \textit{security by design}}
    \item \textbf{Gestione delle operations}
    \item \textbf{Gestione dell'innovazione}
    \item \textbf{Training e \textit{awareness}}
\end{enumerate}
Tale suddivisione dei compiti è coerente con le task del CISO individuate da ENISA in relazione
agli obblighi stabiliti da NIS2\footcite[p. 23]{enisa-cybersecurity-roles}.

\paragraph{Governance e gestione del rischio}
Tra le attività di governance e gestione del rischio rientrano la definizione di
policy di sicurezza e procedure da seguire, l'interazione con gli altri organi aziendali,
le attività di risk assessment, la compliance normativa e la gestione degli audit
nell'ambito della sicurezza informatica.

Il ciclo di risk management richiede infatti un insieme di policy coerente (una principale a livello di business corredata da altre 
specifiche per ogni area di rischio identificata), indicatori/KPI per monitorare efficacia e maturità delle misure, e riesami pianificati o
basati su incidenti significativi, cambiamenti organizzativi o di minaccia, con
decisioni di accettazione del rischio e piani correttivi approvati dal management\footcite[pp. 9--15, pp. 21--25]{enisa-implementing-guidance}.

 \paragraph{\textit{Security by design}}
Nella sfera della \textit{security by design} si trovano tutte le azioni
volte alla messa in sicurezza e al mantenimento della sicurezza degli asset
digitali dell'impresa. L'analisi statica (SAST) e dinamica (DAST) del codice degli
applicativi, le attività di \textit{vulnerability assessment} e \textit{penetration testing}
antecedenti al rilascio dei servizi digitali, le pratiche di \textit{patching} e
aggiornamento periodico sono tutte attività necessarie per prevenire eventuali
incidenti e ridurre così in maniera significativa il rischio residuo.

\paragraph{Operations}
Nell'ambito di \textit{cybersecurity operations} si collocano tutte le procedure
di gestione degli eventi di sicurezza, ossia l'analisi e il monitoraggio degli eventi 
di minaccia, le modalità di risposta, la gestione delle crisi che possono verificarsi
in seguito a un incidente. Tutto ciò che fa parte della difesa attiva e passiva
dell'organizzazione è dunque oggetto di interesse primario per il CISO.

\paragraph{Innovazione}
Come sottolineato da Herin e Rivieri, l'innovazione non è una tematica confinata alle attività
di ricerca e sviluppo dei prodotti e dei processi, ma rappresenta un elemento di
interesse trasversale in tutte le aree aziendali, specialmente per la sicurezza
informatica. L'evoluzione degli strumenti e dei prodotti digitali porta infatti
innumerevoli nuove minacce e potenzialità di attacco che ogni azienda deve prontamente
identificare e arginare quanto prima. Le attività di \textit{cyber threat intelligence},
ricerca e conoscenza delle moderne tecniche a disposizione dei cybercriminali, 
identificazione di nuove vulnerabilità ed exploit, rappresentano la prima attività
a difesa dell'organizzazione contro i fattori esterni.

In quest'ottica, il continuo dialogo con la community rappresenta uno strumento
efficace attraverso cui gestire l'innovazione. Herin e Rivieri partecipano alla community CISO
italiana, una tra le più sviluppate e proattive in questo ambito, la quale rappresenta
una rete di esperti della sicurezza informatica che documentano le nuove scoperte e le
condividono al fine di diffondere la conoscenza e aiutare la protezione delle aziende.

\paragraph{Formazione}
\begin{citazione}{mitnick}
    Il fattore umano è l'anello più debole della sicurezza.
\end{citazione}
La citazione di Kevin Mitnick, hacker tra i più noti della storia, individua
la principale problematica delle aziende quando si parla di sicurezza, in ambito
informatico e non: il fattore umano. Come riportato nell'introduzione all'elaborato, 
il social engineering è tra le pratiche più usate per penetrare
le organizzazioni e le stesse devono quindi investire nella formazione
dei propri dipendenti, oltre a rinforzare i propri sistemi, per evitare accessi
non autorizzati, furti di dati e altri cybercrimini. Qui entra in gioco nuovamente il CISO,
che ha il compito di coordinare la formazione non solo dei dipendenti, ma anche
della direzione aziendale responsabile anche delle decisioni in questo ambito.
Herin sottolinea come la comunicazione in questo processo svolga un ruolo centrale:
è fondamentale adottare una comunicazione adeguata al target a cui si rivolgono le informazioni
per fare in modo che il messaggio sia recepito in maniera chiara e che l'interlocutore
si senta responsabile delle proprie azioni nel quadro della sicurezza aziendale.
Parte del processo di educazione alla sicurezza informatica di competenza del CISO 
riguarda quindi l'individuazione delle aree e dei soggetti più a rischio per 
allocare le risorse formative nel modo più efficace possibile per l'organizzazione.


\subsection{Evoluzione del ruolo del CISO}
La definizione dei compiti del CISO sopra riportata consente di riflettere attentamente
su come è cambiato questa funzione aziendale con la forte evoluzione delle minacce informatiche
e i nuovi approcci delle organizzazioni e delle autorità per gestirle.

Come emergerà anche dalle sezioni successive, il CISO concentra le proprie attività non più esclusivamente sugli aspetti tecnici 
e sull’expertise individuale, ma sulla leadership, sulla gestione del rischio e sull’allineamento con gli obiettivi di business. 
La sua funzione evolve da quella di mero specialista operativo a quella di facilitatore di business, capace di integrare la sicurezza 
informatica nei processi decisionali e strategici dell’organizzazione, traducendo il linguaggio tecnico in termini di rischio, 
valore e priorità per l’impresa\footcite[pp. 32--34]{kouns-ciso}.

Per riassumere visivamente la recente evoluzione di questo ruolo e dei compiti che è chiamato a svolgere, si riporta
di seguito una tabella tratta da \cite{kouns-ciso} che racchiude le principali differenze tra il CISO di ieri e quello di oggi.

\begin{table}[H]
\centering
\renewcommand{\arraystretch}{1.5}
\caption{Evoluzione della descrizione del ruolo del CISO\footcite[Tab. 1]{kouns-ciso}}
\begin{tabular}{p{7cm} p{7cm}}
\toprule
\textbf{CISO di ieri} & \textbf{CISO di oggi} \\
\midrule
Esperto tecnico & Consulente di fiducia (\textit{Trusted adviser}) \\
Analista & Facilitatore e leader (\textit{Facilitator and leader}) \\
Esperto di rischio tecnico & Responsabile della gestione del rischio (\textit{Risk manager}) \\
Contributore individuale & Pensatore strategico integrato nel business (\textit{Integrative business thinker}) \\
Chief Security Officer & Chief Risk Officer \\
Amministratore & Stratega (\textit{Strategist}) \\
Manager & Visionario (\textit{Visionary}) \\
\bottomrule
\end{tabular}
\end{table}



% 5.2 %%%%%%%%%%%%%%%%%%%%%%%%%%%%%%%%%%%%%%%%%%%%%%%%%%%%%%%%%%%%%%%%%%%%%%%%
\section{Con chi interagisce il CISO}
Come visto, il CISO svolge funzioni di carattere trasversale,
interessando molteplici aree aziendali e interagendo con livelli profondamente 
differenti tra loro per competenze, compiti e ruolo gerarchico.
È naturale dunque che il CISO sia quotidianamente a contatto con diversi interlocutori,
i cui principali, secondo Herin, sono:
\begin{itemize}
    \item il dipartimento IT, per l'implementazione delle misure di sicurezza informatica e
    le attività di monitoraggio dei sistemi;
    \item i dipartimenti digital e marketing, che gestiscono ad esempio il sito web e l'applicazione 
    mobile dell'azienda, ove presenti;
    \item l'area di produzione, per quanto riguarda la gestione delle componenti di \textit{industrial 
    technology};
    \item il dipartimento HR, per le attività di formazione;
    \item l'area legale, per la compliance normativa nell'area della cybersicurezza;
    \item l'area di internal audit, al fine di comprendere la struttura e gli asset aziendali e
    l'implementazione di un framework di controllo periodico;
    \item il dipartimento per le pubbliche relazioni, utile nel caso in cui sia 
    necessario comunicare con gli stakeholder esterni all'azienda, ad esempio per la
    comunicazione di incidenti di sicurezza che potrebbero provocare loro danni;
    \item la sicurezza fisica, fondamentale per il controllo dell'accesso fisico alle
    macchine, prima porta di ingresso ai sistemi e ai dati aziendali da proteggere;
\end{itemize}

In questo complesso contesto di interazioni, la guida ENISA sui ruoli per la gestione della cybersicurezza
individua un ulteriore partner naturale del CISO
nella figura del \textit{Cyber Legal, Policy and Compliance Officer}, appartenente all'area legale\footcite[pp.~24--25]{enisa-cybersecurity-roles}
Tale funzione può coadiuvare il CISO nello svolgimento dei propri compiti, soprattutto nel caso di aziende di 
grandi dimensioni operanti su industrie molto differenziate tra loro e conseguentemente soggette a numerosi
requisiti di carattere normativo.


\subsection{Dialogo con gli organi di gestione}
Gli obblighi introdotti dalla nuova Direttiva NIS2 hanno rafforzato la necessità del CISO di dialogare
con le figure dell'area dirigenziale e del top management per comunicare le strategie,
i risultati, i dati e le analisi periodiche sullo stato di salute cibernetica dell'azienda.
Pur essendo le strutture aziendali molto diverse tra loro, è possibile delineare una
metodologia di base comune per il dialogo in materia di sicurezza informatica, in cui
il CISO ricopre il ruolo di interlocutore strategico del CdA\footcite[pp. 25--38]{kouns-ciso}.

Solitamente, spiega Herin, si tiene un comitato di sicurezza periodico ogni
mese o ogni due mesi, in cui il team di cybersecurity espone lo stato dell'arte
e sottomette al comitato stesso le decisioni da prendere. A questo tavolo partecipano 
i direttori delle aree IT, sicurezza fisica, audit, privacy ed eventualmente
referenti per la funzione HR o altre funzioni interessate dalle tematiche trattate.
Si tratta quindi di un comitato tecnico, con l'obiettivo di condividere le analisi
sulla situazione corrente, confrontarsi sul raggiungimento degli obiettivi e 
definire le priorità strategiche per la sicurezza.

Più di rado, circa una o due volte all'anno, si svolge invece il comitato di direzione
per la sicurezza informatica, cui partecipa una delegazione del team di cybersecurity 
e il comitato di direzione aziendale. Si tratta in questo caso di una discussione
ad alto livello, con l'obiettivo di definire il budget, il profilo di rischio
aziendale e i cambiamenti in atto. In questa fase si sottopongono le 
questioni critiche al management, che ha il compito di assumere le adeguate decisioni.

Dalla discussione con Herin emerge anche una metodologia precisa per la costruzione
della strategia e del budgeting: cybersecurity non ha obiettivi stabiliti a priori, 
bensì deve fronteggiare la sfida di adattarsi ai cambiamenti introdotti dalle decisioni strategiche 
nelle aree \textit{core business}, definendo le proprie in relazione alle criticità introdotte.
Si tratta quindi di un approccio adattivo. Herin osserva infatti che l'obiettivo della cybersicurezza 
sia quello di assistere le funzioni principali, permettendo loro di operare efficacemente
ed efficientemente, ma rispettando determinati standard di sicurezza.



% 5.3 %%%%%%%%%%%%%%%%%%%%%%%%%%%%%%%%%%%%%%%%%%%%%%%%%%%%%%%%%%%%%%%%%%%%%%%%
\section{Gestione del rischio}
Herin evidenzia come, ad oggi, le attività di cybersecurity aziendale
siano fortemente legate all'analisi, al trattamento e alle strategie di gestione
del rischio e non più alla mera supervisione dell'implementazione di meccanismi di difesa informatici. 
Il profilo di rischio desiderato dall'azienda guida infatti le 
decisioni strategiche e di investimento, con l'obiettivo di allocare le risorse 
nel modo più efficiente possibile, minimizzando il rischio residuo.

Per il raggiungimento di tale obiettivo le aziende, in particolar modo quelle più 
strutturate, dispongono di un \textit{risk office} che si occupa di:
\begin{itemize}
    \item effettuare periodicamente attività di \textit{risk assessment} per analizzare il 
    profilo di rischio corrente dell'azienda;
    \item elaborare un \textit{risk framework} sulla base delle rilevazioni
    svolte e delle necessità individuate, con l'obiettivo di fornire uno strumento
    per uniformare le decisioni per la gestione del rischio a livello aziendale.
\end{itemize}

Il risk framework aiuta quindi il CISO a valutare il livello di rischio di ogni
servizio digitale\footcite[pp. 40--42]{kouns-ciso}, attraverso la stima di tre fattori:
\begin{enumerate}
    \item minaccia, ovvero una valutazione della fonte di rischio;
    \item probabilità, cioè la facilità con cui una minaccia può concretizzarsi
    per l'azienda nello specifico;
    \item impatto, ossia le conseguenze che la minaccia può avere sull'organizzazione.
\end{enumerate}
Il framework permette di calcolare un valore (secondo il \textit{Common vulnerability
Scoring System}) a partire da questi tre fattori e stabilisce le soglie di
accettabilità del rischio. Seguendo questa indicazione, il CISO decide se:
\begin{itemize}
    \item cessare il rischio;
    \item trattare il rischio;
    \item trasferire il rischio;
    \item accettare il rischio.
\end{itemize}
In particolare, Herin sottolinea come il CISO sia responsabile per le decisioni
in merito ai rischi che riguardano strettamente la cybersicurezza. In linea generale, chi è responsabile 
dei rischi di un asset digitale è l'\textit{application owner} di tale asset, definendo
un preciso framework di \textit{ownership} e responsabilità. Così, ad esempio,
il CISO è responsabile degli applicativi \textit{cyber for cyber} dell'azienda,
ma non per i rischi connessi ad applicativi esterni, i quali avranno un proprio
responsabile definito.



% 5.4 %%%%%%%%%%%%%%%%%%%%%%%%%%%%%%%%%%%%%%%%%%%%%%%%%%%%%%%%%%%%%%%%%%%%%%%%
\section{Adempimento della Direttiva NIS2 nelle organizzazioni}
\subsection{Ruolo e responsabilità del CISO}
In materia di NIS2, il CISO svolge la funzione di punto di contatto richiesta
dal legislatore. La compliance normativa in ambito cyber è infatti uno dei
compiti a carico del CISO, che si occupa di definire le politiche, le procedure
e i comportamenti da adottare nell'organizzazione per conformarsi, in questo caso, alla Direttiva NIS2.

Un primo elemento di novità introdotto da NIS2 nei compiti del CISO, osservano Herin e Rivieri, è il fatto
che la direttiva non fornisca semplicemente controlli da implementare e adempimenti
da dimostrare, ma introduca anche la valutazione del rischio in questo processo.
Infatti, le misure stabilite nella norma, oltre ad applicarsi differentemente a seconda
della classificazione del soggetto in \textit{essenziale} e \textit{importante},
si attuano anche in funzione del rischio. Può accadere quindi che, per due soggetti classificati
ugualmente e operanti in business differenti, alcuni controlli si ritengano
necessari nell'uno e non nell'altro caso. Spetta dunque al CISO, sulla base di un'analisi
di rischio rispetto al profilo richiesto dalla normativa all'azienda, individuare
le vulnerabilità critiche da sanare e presentarle al rispettivo \textit{risk owner}
per decidere come gestire tale fattore di rischio.

Inoltre, in entrambe le interviste con i CISO emerge una questione importante riguardo alle responsabilità in
ambito NIS2. La normativa infatti specifica esplicitamente che i responsabili
dell'implementazione delle misure stabilite siano gli organi dirigenziali dell'azienda,
mentre il CISO non è imputabile direttamente come responsabile dalle autorità
in caso di mancato adempimento. Il suo ruolo è infatti a livello più tecnico e operativo: 
ciò che rientra nelle sue competenze è l'attuazione concreta dei controlli
stabiliti dal framework e l'individuazione delle vulnerabilità, ponendo quindi 
tra le sue responsabilità la corretta realizzazione dei meccanismi di cybersicurezza, 
ma non l'adempimento normativo in senso stretto. In altri termini,
il CISO è la figura deputata a \textit{far funzionare correttamente la macchina} che 
implementa il framework NIS2, sovrintendendone l'adozione delle misure.

Appreso dunque che NIS2 non introduce nessuna nuova responsabilità operativa per la figura
del CISO, Michele Rivieri sottolinea come, invece, abbia introdotto obblighi sul piano informativo
e relazionale: dal momento che il CdA è responsabile della sicurezza, il compito del CISO
è quello di informare periodicamente il top management riguardo i rischi informatici cui 
l'azienda è esposta, il profilo attuale, quello desiderato e le aree di investimento in cui
operare a livello strategico. Tale interazione era presente anche in precedenza, ma Rivieri
precisa che NIS2 ne ha modificato la natura da due punti di vista:
\begin{itemize}
    \item l'attribuzione formale di responsabilità al management implica una maggiore frequenza
        di interazione con gli organi deputati a decidere in merito ai rischi informatici, che devono
        essere informati periodicamente, come descritto nelle sezioni precedenti;
    \item i \textit{risk owner} veri e propri, ovvero coloro che devono decidere in merito alla
        gestione dei rischi dei propri applicativi, fanno capo al CdA in quanto organo che ha capacità
        di spesa in caso di non accettazione del rischio.
\end{itemize}

In questo contesto, la Direttiva NIS2 introduce un significativo cambiamento nei rapporti di forza interni alle organizzazioni: 
l’attribuzione di responsabilità dirette al top management per la conformità alle misure di sicurezza — e per gli eventuali 
inadempimenti — trasforma la cybersicurezza da questione meramente tecnica a tema strategico e di governance.
La prospettiva di sanzioni economiche e reputazionali, unite alla possibilità di responsabilità personali degli amministratori, 
induce il management a una nuova consapevolezza del rischio informatico e ad attribuire maggiore attenzione alle raccomandazioni 
provenienti dall’area di sicurezza. Ne deriva che il CISO, da figura spesso confinata in un ruolo consulenziale e tecnico, assume 
una posizione di effettiva influenza sulle decisioni aziendali, potendo incidere in modo sostanziale sulla definizione delle 
priorità di investimento, sulle strategie di gestione del rischio e sull’integrazione della sicurezza nei processi di business. 
Dal dialogo con i CISO emerge quindi la necessità di formare le figure addette alla cybersicurezza non solo dal punto di vista
tecnico e operativo, ma anche sotto gli aspetti economici, strategici e organizzativi, parte integrante della gestione
della sicurezza in ambito aziendale.


\subsection{Delega di responsabilità}
La gestione della responsabilità secondo NIS2 avviene in maniera rigida e inderogabile: la normativa assegna
la responsabilità della sicurezza informatica agli organi di gestione e per adempiere a tale obbligo non è 
dunque possibile adottare meccanismi di delega.

Il CISO di Prometeia sottolinea come questo obbligo normativo costituisca un problema rilevante
a livello organizzativo, su cui ogni azienda deve interrogarsi e sviluppare un processo ben definito:
se è vero che non è possibile delegare la sicurezza al CISO è allora necessario sottomettere ogni decisione
in ambito di sicurezza al CdA. Tuttavia, come osserva Rivieri, questo assetto comporta un notevole rallentamento 
dei processi decisionali e la mancanza di competenze specifiche di sicurezza informatica determina
un'ulteriore inibizione della capacità decisionale in materia.

Non essendo dunque possibile richiedere agli organi di gestione di farsi carico di tutte le scelte strategiche 
sulla cybersicurezza, come invece un'applicazione ligia di NIS2 richiederebbe, le aziende devono implementare
una qualche forma di delega implicita. La soluzione migliore secondo Rivieri rimane il confronto periodico 
con il CdA, rimandando a tale circostanza le decisioni più impattanti a livello di business, mentre quelle
di rilevanza minore possono essere gestite direttamente dal CISO e dal team cybersecurity. La sicurezza
informatica è infatti un processo giornaliero, data la grande velocità con cui evolvono gli scenari, e 
rinunciare al meccanismo della delega implicherebbe uno svantaggio notevole nella gestione delle minacce cyber.

Tuttavia, Herin e Rivieri concordano nel riconoscere un ruolo centrale alla formazione del top management:
senza la possibilità di comprendere i rischi che l'azienda affronta gli organi di gestione non hanno
capacità decisionale e la definizione di una strategia di cybersicurezza diventerebbe la mera accettazione
di quanto presentato dal CISO. Ciò andrebbe contro la normativa NIS2, che riconosce nel CdA un ulteriore
organo di controllo all'interno della catena della sicurezza aziendale, con un ruolo attivo 
nella gestione dei rischi che l'impresa affronta. 



% 5.5 %%%%%%%%%%%%%%%%%%%%%%%%%%%%%%%%%%%%%%%%%%%%%%%%%%%%%%%%%%%%%%%%%%%%%%%%
\section{Ruolo degli standard secondo i CISO}
La struttura di NIS2, dettagliata nella descrizione dei controlli e delle responsabilità,
ne permette un utilizzo diretto proprio come un framework di cybersecurity.
Herin sostiene infatti che, laddove l'azienda non abbia già implementato standard
per la cybersicurezza, l'approccio migliore sia quello di progettare la strategia
di sicurezza informatica direttamente sulla normativa. In questo modo, le aziende 
possono soddisfare gli obblighi imposti dal regolatore in maniera mirata, 
strutturando la governance e le policy in modo da essere conformi a quanto richiesto. 

Piuttosto che adottare uno standard per favorire l'applicazione di NIS2, Herin suggerisce
il processo contrario: visti i numerosi punti in comune, una volta fissate le misure
per NIS2 è possibile anche adattare le scelte in funzione dell'adozione di ISO 27001 o del NIST CSF 2.0,
piuttosto che di altri standard che l'azienda individua come strategicamente rilevanti.
Infatti, se passare da uno standard per arrivare alla normativa può apparire controintuitivo,
data l'obbligatorietà della seconda e la natura volontaria del primo, il processo contrario
può portare diversi vantaggi: da un lato, consente di concentrare le risorse sulla realizzazione dei
requisiti obbligatori, garantendo la compliance, e dall'altro struttura tale base per la futura
adozione di framework, i quali ricicleranno i controlli compatibili già implementati, 
aggiungendone di nuovi o adattando quelli esistenti ove necessario.

Tuttavia, è importante notare che, trattandosi di una normativa e non un testo tecnico, il linguaggio
adottato e le modalità di presentazione dei controlli lasciano un certo spazio all'interpretazione\footcite[p. 3]{parmar-miles-csf}.
Ciò dipende in larga misura anche dalla maturità delle procedure di sicurezza informatica già adottate dall'azienda: 
i soggetti con un livello di maturità della sicurezza superiore avranno la necessità di individuare solamente 
le aree in cui risultano carenti rispetto a NIS2, mentre le entità che per la prima volta rientrano in un 
regime di regolamentazione delle misure di cybersicurezza dovranno implementare una strategia efficace da zero.
In quest'ultima casistica, che può presentare notevoli sfide per le aziende, come visto nel Capitolo \ref{cap:nis2},
gli standard di cybersicurezza rappresentano un valido strumento a supporto dell'organizzazione e del CISO. 
Per arrivare al piano d'azione concreto a partire dalla normativa, la letteratura suggerisce infatti di elaborare una
strategia integrando le misure NIS2 nel contesto di uno (o più) framework di sicurezza informatica, secondo lo schema
presentato in figura (Fig. \ref{fig:implementazione-nis2})\footcite[p. 5]{parmar-miles-csf}.

\defaultimage{impl-NIS2}{Flusso di implementazione della direttiva NIS2\footcite[Fig. 1]{parmar-miles-csf}}{fig:implementazione-nis2}

I framework, dunque, si inseriscono nel contesto organizzativo come facilitatori per l'adozione di NIS2, creando una solida
base di controlli, procedure e governance da integrare secondo gli obblighi specifici previsti dalle normative,
ma non essenziali per l'adempimento di NIS2, come evidenziato dall'analisi condotta nel Capitolo \ref{cap:confronto}.
